% SPDX-License-Identifier: LPPL-1.3c
% Copyright (C) 2026 Christophe Jorssen
% Author and Current Maintainer: Christophe Jorssen <christophe.jorssen@gmail.com>
% LPPL maintenance status: maintained
% This file is part of luafemm. See LICENSE and LICENSES.md for its terms.
\documentclass[a4paper,11pt]{article}
\usepackage[margin=19mm]{geometry}
\usepackage{fontspec}
\usepackage{amsmath}
\usepackage{luafemm}
\pagestyle{empty}
\definecolor{field}{HTML}{126E82}
\definecolor{copper}{HTML}{CA783E}
\begin{document}
\noindent{\LARGE\bfseries An electromagnet computed in LuaTeX}\par
\smallskip
\noindent\textsf{luafemm \luafemmversion{} · nonlinear planar magnetostatics}

\medskip
\noindent TikZ paths describe the materials and draw the device. Lua solves
for the finite-element vector potential; TikZ overlays oriented field lines.

\begin{center}
\tikzset{
  % Illustrative B (T) / H (A/m) curve, not a commercial steel grade.
  femm/materials/iron/.style={bh={0/0,0.5/80,1.0/180,1.3/350,
    1.5/900,1.65/2500,1.8/7000,2.0/20000}},
  core/.style={femm/region={material=iron},draw=black!65,
    fill=black!9,line width=.65pt},
  winding/.style={femm/region={material=copper},draw=copper!80!black,
    fill=copper!25,line width=.65pt}
}
\begin{tikzpicture}[
  femm/problem={xmin=-110,xmax=110,ymin=-95,ymax=115,
    mesh size=2,mesher=delaunay,curve tolerance=.03}]
  \filldraw[core] (-40,-30)--(40,-30)--(40,30)--(25,30)
    --(25,-15)--(-25,-15)--(-25,30)--(-40,30)--cycle;
  \filldraw[core] (-40,33) rectangle (40,48);
  % NI = 6000 ampere-turns; each winding section measures 6 x 26 square mm.
  \filldraw[winding,femm/current density=6000/(6*26)*1000000]
    (-48,-10) rectangle (-42,16);
  \filldraw[winding,femm/current density=-6000/(6*26)*1000000]
    (-23,-10) rectangle (-17,16);
  \begin{scope}
    \clip (-64,-47) rectangle (65,66);
    \pic[femm/lines=29,femm/field style={color=field}] {femm field};
  \end{scope}
  % Draw conductor symbols after the field to keep them readable.
  \foreach \y in {-6,0,6,12}{
    \node[fill=copper!25,inner sep=.2pt,text=copper!65!black] at (-45,\y) {$\odot$};
    \node[fill=copper!25,inner sep=.2pt,text=copper!65!black] at (-20,\y) {$\otimes$};
  }
  \node[fill=white,inner sep=2pt,font=\small] at (0,40.5) {Armature};
  \node[fill=white,inner sep=2pt,font=\small] at (0,-23) {U-shaped core};
  \node[align=center,font=\small] at (0,7) {Air\\$NI=6000\,\mathrm{A\!\cdot\!turns}$};
  \draw[thin] (40,30)--(52,30) (40,33)--(52,33);
  \draw[thin,-Stealth] (49,25)--(49,30);
  \draw[thin,-Stealth] (49,38)--(49,33);
  \node[anchor=west,font=\small,fill=white,inner sep=1pt] at (51,31.5) {$e=3\,\mathrm{mm}$};
  \draw[thin] (-40,-30)--(-40,-39) (40,-30)--(40,-39);
  \draw[Stealth-Stealth,thin] (-40,-37)--(40,-37)
    node[midway,fill=white,inner sep=2pt,font=\small] {$80\,\mathrm{mm}$};
  \node[anchor=west,font=\small,text=field,fill=white,inner sep=2pt] at (-60,61) {\bfseries Contours $A_z=\mathrm{constant}$, oriented along $\mathbf B$};
  \node[anchor=west,font=\small,fill=white,inner sep=2pt] at (-60,-44) {$\odot\ J_z>0$ (outward)\qquad $\otimes\ J_z<0$ (inward)};
\end{tikzpicture}
\end{center}
% The solved model stays available in Lua after the TeX group closes.
\noindent\textbf{Computed result.}
At the gap centres, $|\mathbf B|=\femmB{-32.5}{31.5}\,\mathrm T$
on the left and $\femmB{32.5}{31.5}\,\mathrm T$ on the right.
Mesh: \femmstat{nodes} nodes, \femmstat{triangles} triangles.
Newton: \femmstat{newton} iterations.

\medskip
\noindent\textbf{Model.} The ferromagnetic section is $15\,\mathrm{mm}$ wide,
with two $3\,\mathrm{mm}$ gaps and infinite depth.
The interpolated $H(B)$ curve describes an illustrative saturating material.
The outer boundary of the $220\times210\,\mathrm{mm}^2$ air domain has
$A_z=0$; the figure shows a closer crop.

\medskip
\noindent\textbf{Computation.}
$\mathbf B=(\partial_y A_z,-\partial_x A_z)$ and
$-\nabla\!\cdot\!\bigl(\nu(|\nabla A_z|)\nabla A_z\bigr)=J_z$,
with $\nu=H/B$.
The solver uses triangular $P_1$ elements, damped Newton iterations and
preconditioned conjugate gradients. The contours follow the computed
potential without graphical smoothing.
\end{document}
